\begin{activity}{Thermal vs Chemical Diffusion}

\section*{Purpose}
To compare transient thermal diffusion with chemical diffusion processes familiar from igneous and metamorphic petrology.

\section*{Diffusion Framework}
Both thermal and chemical diffusion can be written in the general form
\[
\frac{\partial \phi}{\partial t} = D \nabla^2 \phi,
\]
where $\phi$ represents temperature or chemical concentration, and $D$ is a diffusivity.

\section*{Tasks}
\begin{enumerate}
  \item Explain why both thermal and chemical diffusion act to smooth spatial gradients over time.

  \item Chemical diffusion in minerals often preserves sharp zoning. Give two reasons why sharp compositional gradients may survive despite diffusion.

  \item Describe how multi-component diffusion differs from thermal diffusion involving a single scalar field.
\end{enumerate}

\section*{Comparison}
\begin{itemize}
  \item Thermal diffusion: single field, relatively uniform diffusivity.
  \item Chemical diffusion: multi-component, potentially rate-limited by the slowest species.
\end{itemize}

\section*{Reflection}
Why is heat generally a poorer long-term recorder of geological processes than composition?

\end{activity}